%##################################################################################
 
\documentclass[12pt,a4paper,twoside]{report}
%\documentclass[12pt,a4paper,oneside]{report}
\usepackage{graphicx}
\usepackage[english]{babel} %-- uncomment this to get english titles
\usepackage{german,a4}
%\usepackage{picins}
%\usepackage{epsfig}
\usepackage{fancyhdr}			% for nice header and footer
\usepackage{hyperref}			% references in pdf
\usepackage[sf]{titlesec}	%customization of \chapter Titles for appendices
\usepackage{textcomp}
\usepackage{makeidx}
\usepackage{ngerman}			% german Umlauts
\usepackage[utf8]{inputenc}
\usepackage{booktabs}                               % necessary for tabulars
\usepackage[squaren]{SIunits}
\usepackage[numbers,square]{natbib}
 
%\usepackage{wrapfig}
%\usepackage[pdftex]{graphicx}
%\usepackage{ifthen}  
%\usepackage{booktabs} % fancy tables
%\usepackage[titletoc]{appendix} % custom naming of appendices
\usepackage{amssymb, amsmath, amsthm} % for equations & eqref
%\usepackage{listings} \lstset{basicstyle=\tiny\ttfamily, numbers=left, escapeinside={(*}{*)}, captionpos=b}
%\usepackage{dirtree}

%get bigger \par with one empty line
\newcommand{\mypar}{\par\medskip}

%TODO line
\newcommand{\writeTodo}{1}
\newcommand{\todo}[1]{
	\ifdefined \writeTodo
		\mypar\textbf{\textcolor{KITgreen}{TODO: }\textcolor{red}{#1}}\mypar
	\fi
} 
 
%\theoremstyle{plain}% default
\theoremstyle{definition}
\newtheorem{definition}{Definition}

% Abbkürzungsverzeichnis einfügen
%%%%%%%%%%%%%%%%%%%%%%%%%%%%%%%%%%%%%%%%%%%%%%%%%%%%%%%%% 
\usepackage{nomencl}
\let\abbrev\nomenclature
\renewcommand{\nomname}{Abkürzungsverzeichnis}
\setlength{\nomlabelwidth}{.25\hsize}
\renewcommand{\nomlabel}[1]{#1 }%\dotfill}
\setlength{\nomitemsep}{-\parsep}
\makenomenclature
\newcommand{\markup}[1]{\underline{#1}} 

% Farben
\usepackage{color}
\definecolor{KITgreen}{rgb}{0, .61, .51} 
\definecolor{KITbluegrey}{rgb}{.27, .39, .67} 
\definecolor{KITgrey}{rgb}{.49, .49, .49} 

% =====================================================
% Dokumenten-Platzhalter
% =====================================================
% =====================================================
% Dokumenten Platzhalter
% =====================================================


\newcommand{\titelderarbeit}{TalkyCars: A Distributed Software Platform for Cooperative Perception among Connected Autonomous Vehicles based on Cellular-V2X Communication}
\newcommand{\artderarbeit}{Master Thesis}

\newcommand{\diplomandprefix}{B. Sc.}
\newcommand{\diplomand}{Ferdinand Mütsch}

\newcommand{\betreuerA}{Marco Stang}
\newcommand{\betreuerB}{Martin Böhme}

 
%\newcommand{\nameprefix}{}{}    %use no nameprefix 
\newcommand{\nameprefix}{M. Sc. }
\newcommand{\docauthor}{Andreas Lauber}

\newcommand{\nameprefixb}{M. Sc. }    %use no nameprefix 
\newcommand{\docauthorb}{}

\newcommand{\nameprefixc}{}{}    %use no nameprefix 
\newcommand{\docauthorc}{}

%\newcommand{\betreuerA}{nn}
%\newcommand{\betreuerB}{nn}
\newcommand{\abgabe}{1 August 2016}
\newcommand{\versionierungsnr}{v1.0}
\newcommand{\fussnoteninhalt}{}
\newcommand{\leerzeichen}{ }

%% english version
%\newcommand{\Dachorganisation}{Karlsruhe Institute of Technology - KIT}
%\newcommand{\Institut}{Institute for Information Processing 
%												Technology - ITIV}
												
% deutsche version
\newcommand{\Dachorganisation}{Karlsruher Institut für Technologie - KIT}
\newcommand{\Institut}{Institut für Technik der Informationsverarbeitung - ITIV}

% ====================================================
% Formatierung des Dokuments
% ====================================================


%\usepackage{color}
%\usepackage{helvet}

% =====================================================
% Seite einrichten
% =====================================================
\setlength{\textwidth}{160mm} \setlength{\textheight}{235mm}
\setlength{\oddsidemargin}{5mm}
%pagelayout: (1in=25,4mm), default offset of page origin: 1in x 1in
%     25,4+5=30,4mm       160mm          19,6mm
%  |                   sadfasdfasdfasdf          |
\setlength{\evensidemargin}{-5,8mm} 
%     25,4-5,8= 19,6mm    160mm          30,4mm
%  |                   sadfasdfasdfasdf          |
\setlength{\topmargin}{0mm} \setlength{\voffset}{-15mm}
\setlength{\headsep}{12mm} \setlength{\footskip}{15mm}
\setlength{\headheight}{15.5pt}

% =====================================================
% sonstige Einstellungen
% =====================================================
%\flushbottom % Textfluss schoen unten ausrichten
\footnotesep12pt % Abstand Text / Fussnote
%\setlength{\parindent}{0pt} \setlength{\abovecaptionskip}{-6pt}
%\setlength{\belowcaptionskip}{0pt} \setlength{\intextsep}{18pt}
%\renewcommand{\baselinestretch}{1.2} % Zeilenabstand
\setcounter{secnumdepth}{4}
\newcommand{\p}[1]{\texttt{#1}}
\nonfrenchspacing

% =====================================================
% Kopf- und Fusszeile formatieren
% =====================================================
\pagestyle{fancy}
\fancyhead{}
\makeatletter
\if@twoside
	\fancyhead[LO]{\sf{\rightmark}}
	\fancyhead[RE]{\sf{\leftmark}}
	\fancyfoot[LE]{\setlength{\unitlength}{1mm}
		\begin{picture}(0,0) \put(0,-2){
			%\includegraphics[height=0.6cm,width=0.6cm]{98_images/ITIVlogo.png}
		} \end{picture}
		%\put(9,2){\scriptsize\sf{\Institut}}\put(9,-2){\scriptsize\sf{\Dachorganisation}}}
		\put(0,2){\scriptsize\sf{\Institut}}\put(0,-2){\scriptsize\sf{\Dachorganisation}}}

\else
	\fancyhead[LO]{\sf{\leftmark}}
	\fancyfoot[LO]{\setlength{\unitlength}{1mm} 
		\begin{picture}(0,0) \put(0,-2){
			% \includegraphics[height=0.6cm,width=0.6cm]{98_images/ITIVlogo.png}
		} \end{picture}
		%\put(9,2){\scriptsize\sf{\Institut}}\put(9,-2){\scriptsize\sf{\Dachorganisation}}}
		\put(0,2){\scriptsize\sf{\Institut}}\put(0,-2){\scriptsize\sf{\Dachorganisation}}}
\fi
\fancyhead[LE, RO]{\sf{\thepage}}
\makeatother

\cfoot{}
\fancyfoot[RO]{\footnotesize\sf{\titelderarbeit}}
\renewcommand{\headrulewidth}{0.4pt}
\renewcommand{\footrulewidth}{1.0pt}

% ====================================================
% Kopf- und Fusszeile fuer "plain"-Format ueberschreiben
% ====================================================
\makeatletter
\if@twoside
	\fancypagestyle{plain}{%
	\fancyhf{}
	\fancyfoot[LE]{\setlength{\unitlength}{1mm}
		\begin{picture}(0,0) \put(0,-2){
			%\includegraphics[height=0.6cm,width=0.6cm]{98_images/ITIVlogo.png}
		} \end{picture}
		%\put(9,2){\scriptsize\sf{\Institut}}\put(9,-2){\scriptsize\sf{\Dachorganisation}}}
		\put(0,2){\scriptsize\sf{\Institut}}\put(0,-2){\scriptsize\sf{\Dachorganisation}}}
	\cfoot{}
	\fancyfoot[RO]{\footnotesize\sf{\titelderarbeit}}
	\renewcommand{\headrulewidth}{0pt}
	\renewcommand{\footrulewidth}{1.0pt}}
\else
	\fancypagestyle{plain}{%
	\fancyhf{}
	\fancyfoot[L,C,R]{\setlength{\unitlength}{1mm}
	\begin{picture}(0,0) \put(0,-2){
		%\includegraphics[height=0.6cm,width=0.6cm]{98_images/ITIVlogo.png}
	} \end{picture}
	%\put(9,2){\scriptsize\sf{\Institut}}\put(9,-2){\scriptsize\sf{\Dachorganisation}}}
	\put(0,2){\scriptsize\sf{\Institut}}\put(0,-2){\scriptsize\sf{\Dachorganisation}}}
	\cfoot{}
	\rfoot{\footnotesize\sf{\titelderarbeit}}
	\renewcommand{\headrulewidth}{0pt}
	\renewcommand{\footrulewidth}{1.0pt}}
\fi

%makes the index do not remove!
\makeindex

\hypersetup{
	%pdftex %schon gesetzt
	pdfauthor = \docauthor,
	%pagebackref,%schon gesetzt
	colorlinks=true,
	linkcolor=black,
	citecolor=black,
	urlcolor=black
	}

\providecommand{\e}[1]{\ensuremath{\times 10^{#1}}}

% =====================================================
% Inhalt der Titelseite definieren
% =====================================================

\makeindex
\newcommand{\Idx}[1]{#1 \index{#1}}

\hyphenation{
EVITA
}

% =====================================================
% Zeichen für Copyright, Trademark, Registerd, ...
% =====================================================
\def\TReg{\textsuperscript{\textregistered}}
\def\TCop{\textsuperscript{\textcopyright}}
\def\TTra{\textsuperscript{\texttrademark}}

% =====================================================
% selbs definierte Zeichen
% =====================================================
\def\zB{z.\,B. }
\def\uvm{u.\,v.\,m.}
\begin{document}

% =====================================================
% Put Titel in English
% ===================================================== 
%
% Inhalt der Titelseite definieren
\pagestyle{empty}
\pagenumbering{roman}
\setcounter{page}{1}

\begin{titlepage}

\begin{minipage}{15cm}
\linespread{1.2}

\begin{tabular}{lr}
\begin{minipage}{0.2\linewidth}
\includegraphics[height=1cm]{98_images/kit.png}
\end{minipage}
&
\begin{minipage}{0.8\linewidth}
\large
%\center
\vspace{7mm}
\begin{flushleft}
\Dachorganisation
\end{flushleft}
\end{minipage}
\\
\\
\begin{minipage}{0.2\linewidth}
\includegraphics[height=1cm]{98_images/ITIVLogo.png}
\end{minipage}
&
\begin{minipage}{0.8\linewidth}
\large
%\center
\vspace{7mm}
\begin{flushleft}
\Institut\\
\end{flushleft}
\end{minipage}
\\
\end{tabular}

\vspace{3cm}

\begin{center}
\Huge
\bfseries  \titelderarbeit
\end{center}

\vspace{1cm}
\begin{center}
\large 
\subtitelderarbeit 
\vspace{1cm}
%\Large
%\textbf{\diplomandprefix\ \diplomand}

\vspace{1cm}
%\Large
%%Berichtsbezeichnung und ID Nummer
 Version 
\versionierungsnr\\
\abgabe%\textbf{\abgabe}
\end{center}
%
\vspace{2.5cm}

\begin{tabular}{rcl}
\bfseries Head of Institute:
&& Prof. Dr.-Ing. J. Becker\\
&& Prof. Dr. rer. nat  W. Stork\\
&& Prof. Dr.-Ing. Eric Sax\\
 \\
	\bfseries Author:         && \nameprefix \leerzeichen \docauthor \\
											      && \nameprefixb \leerzeichen \docauthorb\\
														&& \nameprefixc \leerzeichen \docauthorc \\
											      \end{tabular}											
											      
\end{minipage}
\end{titlepage}
\pagestyle{fancy}
%    \parindent=0pt
%    %\sloppypar
%    \linespread{1.2}
%    \thispagestyle{plain}
%    %\frontmatter
%    %\maketitle

% =====================================================
% Put Titel in Deutsch
% ===================================================== 

% Inhalt der Titelseite definieren
\pagestyle{empty}
\pagenumbering{roman}
\setcounter{page}{1}

\begin{titlepage}

\begin{minipage}{15cm}
\linespread{1.2}

\begin{tabular}{lr}
\begin{minipage}{0.3\linewidth}
\includegraphics[height=2cm]{98_images/kit.png}
\end{minipage}
&
\begin{minipage}{0.7\linewidth}
\large
\center
Karlsruhe Institute of Technology\\
Fakultät für Elektrotechnik
\end{minipage}
\\
\\
\begin{minipage}{0.3\linewidth}
\includegraphics[height=2cm]{98_images/ITIVlogo.png}
\end{minipage}
&
\begin{minipage}{0.7\linewidth}
\large
\center
Institut für Technik der Informationsverarbeitung (ITIV)\\
\end{minipage}
\\
\end{tabular}

% #################################################################
% ### If the Title is very long reduce vspace to less than 2cm
% #################################################################
\vspace{1.5cm}

% #################################################################

\begin{center}
\Huge
\bfseries  \titelderarbeit
\end{center}

\vspace{1cm}
\begin{center}
\large 
\artderarbeit ~by\\
\vspace{1cm}
\Large
\textbf{\diplomandprefix\ \diplomand}

\vspace{1cm}
\Large
\textbf{\abgabe}
\end{center}
%
\vspace{2.5cm}

\begin{tabular}{rcl}
\bfseries Institutsleitung: 
&& Prof. Dr.-Ing. Dr. h.c. J. Becker\\
&& Prof. Dr.-Ing. Eric Sax\\
&& Prof. Dr. rer. nat  W. Stork\\

 \\
	\bfseries Betreuer:       &&\nameprefix  \betreuerA \\
											      && \nameprefixb  \betreuerB \\
\end{tabular}											
\end{minipage}
\end{titlepage}
    \parindent=0pt
    %\sloppypar
    \linespread{1.2}
    \thispagestyle{plain}
    %\frontmatter
    %\maketitle
    \cleardoublepage
    
% =====================================================
% Abstract
% =====================================================     
	\begin{abstract}
		Hier kommt die Zusammenfassung ...


\cleardoublepage

\chapter*{Urheberrecht}

// TODO
	\end{abstract}
	\cleardoublepage
% =====================================================
% Signaturepage
% ===================================================== 
    \chapter*{}
\begin{flushleft}
\vspace{11cm}
Erklärung\\[1cm]
Ich versichere hiermit, dass ich meine Bachelorarbeit (bzw. Masterarbeit) selbständig und unter Beachtung der Regeln zur Sicherung guter wissenschaftlicher Praxis im Karlsruher Institut für Technologie (KIT) in der aktuellen Fassung angefertigt habe. \\
Ich habe keine anderen als die angegebenen Quellen und Hilfsmittel benutzt und wörtlich oder inhaltlich übernommene Stellen als solche kenntlich gemacht.\\[1cm]

Karlsruhe, den \abgabe\\[1cm]
\end{flushleft}

\begin{center}
---------------------------------\\
\diplomand
\end{center}

	\cleardoublepage
% =====================================================
% TOC
% ===================================================== 
    \tableofcontents
    %\clearpage
    \cleardoublepage

% =====================================================
% Main Chapters
% ===================================================== 
    %\mainmatter
    \pagenumbering{arabic}
    \setcounter{page}{1}
    \pagestyle{fancy}
 \normalsize

    
    \chapter{Introduction}
    \label{ch:Introduction}    
    	\chapter{Source Code}
\label{appendix:source_code}

\section{SQL Queries for Traffic Volume Estimation}
\label{sec:appendix:source_code:traffic_volume}

\subsection{Query: Geographic Area}
\inputminted{sql}{97_listings/traffic_volume_1.sql}

\subsection{Query: Total Road Length}
\inputminted{sql}{97_listings/traffic_volume_2.sql}

\subsection{Query: Average Number of Lanes}
\inputminted{sql}{97_listings/traffic_volume_3.sql}
    \cleardoublepage
    
    \chapter{Background}
    \label{ch:Background}
		\chapter{Source Code}
\label{appendix:source_code}

\section{SQL Queries for Traffic Volume Estimation}
\label{sec:appendix:source_code:traffic_volume}

\subsection{Query: Geographic Area}
\inputminted{sql}{97_listings/traffic_volume_1.sql}

\subsection{Query: Total Road Length}
\inputminted{sql}{97_listings/traffic_volume_2.sql}

\subsection{Query: Average Number of Lanes}
\inputminted{sql}{97_listings/traffic_volume_3.sql}
    \cleardoublepage

	\chapter{Related Work}
	\label{ch:Related_Work}  	
		\chapter{Source Code}
\label{appendix:source_code}

\section{SQL Queries for Traffic Volume Estimation}
\label{sec:appendix:source_code:traffic_volume}

\subsection{Query: Geographic Area}
\inputminted{sql}{97_listings/traffic_volume_1.sql}

\subsection{Query: Total Road Length}
\inputminted{sql}{97_listings/traffic_volume_2.sql}

\subsection{Query: Average Number of Lanes}
\inputminted{sql}{97_listings/traffic_volume_3.sql}
    \cleardoublepage    
    
    \chapter{Problem Description}
    \label{ch:Problem_Description}   
   		\chapter{Source Code}
\label{appendix:source_code}

\section{SQL Queries for Traffic Volume Estimation}
\label{sec:appendix:source_code:traffic_volume}

\subsection{Query: Geographic Area}
\inputminted{sql}{97_listings/traffic_volume_1.sql}

\subsection{Query: Total Road Length}
\inputminted{sql}{97_listings/traffic_volume_2.sql}

\subsection{Query: Average Number of Lanes}
\inputminted{sql}{97_listings/traffic_volume_3.sql}
   		
    \chapter{Design \& Implementation}
	\label{ch:Design_Implementaion}  
		\chapter{Source Code}
\label{appendix:source_code}

\section{SQL Queries for Traffic Volume Estimation}
\label{sec:appendix:source_code:traffic_volume}

\subsection{Query: Geographic Area}
\inputminted{sql}{97_listings/traffic_volume_1.sql}

\subsection{Query: Total Road Length}
\inputminted{sql}{97_listings/traffic_volume_2.sql}

\subsection{Query: Average Number of Lanes}
\inputminted{sql}{97_listings/traffic_volume_3.sql}
    \cleardoublepage    

	\chapter{Evaluation}
	\label{ch:Evaluation}  
	\chapter{Source Code}
\label{appendix:source_code}

\section{SQL Queries for Traffic Volume Estimation}
\label{sec:appendix:source_code:traffic_volume}

\subsection{Query: Geographic Area}
\inputminted{sql}{97_listings/traffic_volume_1.sql}

\subsection{Query: Total Road Length}
\inputminted{sql}{97_listings/traffic_volume_2.sql}

\subsection{Query: Average Number of Lanes}
\inputminted{sql}{97_listings/traffic_volume_3.sql}
	\cleardoublepage  
	
	\chapter{Conclusion \& Future Work}
	\label{ch:Conclusion_Future_Work}  
	\chapter{Source Code}
\label{appendix:source_code}

\section{SQL Queries for Traffic Volume Estimation}
\label{sec:appendix:source_code:traffic_volume}

\subsection{Query: Geographic Area}
\inputminted{sql}{97_listings/traffic_volume_1.sql}

\subsection{Query: Total Road Length}
\inputminted{sql}{97_listings/traffic_volume_2.sql}

\subsection{Query: Average Number of Lanes}
\inputminted{sql}{97_listings/traffic_volume_3.sql}
	\cleardoublepage  

    
% =====================================================
% Bibliography
% ===================================================== 
\nocite{*}
%\addcontentsline{toc}{chapter}{Literaturverzeichnis, in ToC löschen}
%\addcontentsline{toc}{chapter}{Literature}
%\bibliographystyle{dinat}        % use a DIN style for the bibliography
%\bibliographystyle{plain}        % use a DIN style for the bibliography
%\bibliographystyle{abbrv}			%mit nummern in []
%\bibliographystyle{plain}			%mit nummern in []
%\bibliographystyle{alpha}			%3 Buchstaben + Jahr
\bibliographystyle{alphadin}

    % use external Bib-File, 
		\bibliography{99_bib/standardbib}		
		\cleardoublepage
% =====================================================
% Appendix
% ===================================================== 
    \appendix

%    \input{95_anhang/anhang}
		\cleardoublepage

    \listoftables
    	\cleardoublepage
    \listoffigures
    	\cleardoublepage
    \printindex
    	\cleardoublepage
    \printnomenclature
    \renewcommand{\leftmark}{\uppercase{Abkürzungsverzeichnis}}
    	\cleardoublepage

\end{document}

